% Analise complementar: resultados exclusivamente por coleta.
% Para a metrica por leitura indicada pelo orientador, usar
% validacao_por_leitura/secao_para_artigo.tex como secao principal.
% Fontes: validacao_cruzada_126_combinacoes.csv e
% exploracao_teste_126_combinacoes.csv.
\subsection{Análise combinatória dos sensores por coleta}

Foram avaliados os 63 subconjuntos não vazios dos seis sensores MQ em dois
cenários: sinais MQ brutos isolados e os mesmos sinais acrescidos do índice de
umidade do solo, temperatura e pressão. Empregou-se Extra Trees com 900 árvores
e hiperparâmetros fixos. As 24 coletas de treinamento foram divididas em quatro
partições estratificadas por classe, sempre mantendo cada coleta integralmente
em uma partição. As probabilidades das leituras de uma coleta foram agregadas
pela média, produzindo uma única previsão por coleta. Todas as métricas e
seleções desta análise foram calculadas nesse nível.

\begin{table}[htbp]
\centering
\caption{Acurácia balanceada por coleta na análise combinatória. CV:
validação cruzada nas 24 coletas de treinamento; teste: 12 coletas.}
\label{tab:combinatoria_mq_ambiente_coleta}
\begin{tabular}{llrrr}
\hline
Cenário & Sensores MQ & $n_{MQ}$ & CV (\%) & Teste (\%) \\
\hline
MQ & MQ7 (melhor CV) & 1 & 71,68 & 75,00 \\
MQ & MQ2, MQ3, MQ8, MQ135, MQ138 (exploratório) & 5 & 66,43 & 91,67 \\
MQ & Todos os seis & 6 & 65,73 & 91,67 \\
MQ + ambiente & MQ2 (melhor CV) & 1 & 74,83 & 83,33 \\
MQ + ambiente & MQ2, MQ135, MQ138 (exploratório) & 3 & 74,13 & 91,67 \\
MQ + ambiente & Todos os seis & 6 & 70,28 & 83,33 \\
\hline
\end{tabular}
\end{table}

Na validação cruzada, a inclusão das variáveis ambientais elevou a acurácia
balanceada por coleta em 59 dos 63 pares de subconjuntos MQ idênticos, reduziu
em um e não a alterou em três. Os melhores subconjuntos selecionados apenas no
treinamento foram MQ7 sem ambiente e MQ2 com ambiente. No teste, acertaram,
respectivamente, 9 e 10 das 12 coletas. Os seis MQ sem ambiente acertaram
11 das 12 coletas, enquanto os seis MQ com ambiente acertaram 10 das 12.

Em uma busca exploratória que avaliou as 126 configurações nas 12 coletas de
teste, 34 empataram no maior desempenho, com 11 acertos. Entre elas, o conjunto
MQ2, MQ135 e MQ138 acrescido de solo, temperatura e pressão apresentou a maior
acurácia balanceada por coleta na validação cruzada do treino (74,13\%).
Restrita aos MQ, a mesma regra selecionou MQ2, MQ3, MQ8, MQ135 e MQ138
(66,43\% na validação). Como o teste participou dessa seleção, 91,67\% descreve
o desempenho observado nessas coletas e não estima de forma imparcial o
desempenho futuro dos conjuntos escolhidos. Novas coletas independentes são
necessárias para confirmar a configuração antes de fixá-la para o hardware.
